% TOP_PROBLEM: {"id": 20001112, "title": "Simultaneously avoiding bounded weighted averages", "category_id": 2, "category": {"id": 2, "name": "combinatorics", "display_name": "Combinatorics", "description": "Counting problems, graph theory, discrete structures.", "slug": "combinatorics", "order_index": 2, "created_at": "2026-07-31T15:26:25.671Z"}, "status": "partially_solved", "problem_number": "AIM-COMBINATORICS-0237", "set_id": 14}
% FIRST_POSTED: 2026-10-01
% ENTRY_KIND: statement_only
% UnsolvedMath ID 20001112; source catalogue code AIM-COMBINATORICS-0237
% !TeX program = lualatex
% Definition and sources only; this file is not an LLM solution attempt.
\documentclass[11pt,a4paper]{article}
\usepackage[margin=25mm]{geometry}
\usepackage{fontspec}
\setmainfont{lmroman10-regular.otf}[BoldFont=lmroman10-bold.otf,
 ItalicFont=lmroman10-italic.otf,BoldItalicFont=lmroman10-bolditalic.otf]
\usepackage{amsmath,amssymb,mathtools}
\usepackage{unicode-math}
\setmathfont{latinmodern-math.otf}
\AtBeginDocument{\let\setminus\smallsetminus}
\usepackage{xurl,hyperref}
\hypersetup{colorlinks=true,urlcolor=blue}
\setcounter{secnumdepth}{0}
\setlength{\parindent}{0pt}
\setlength{\parskip}{5pt}
\setlength{\emergencystretch}{3em}
\begin{document}
\section*{Simultaneously avoiding bounded weighted averages}\label{20001112}
{\small UnsolvedMath ID 20001112 · AIM-COMBINATORICS-0237 · Combinatorics · AIM Workshop Problem Lists}
\subsection{Definitions and mathematical statement}
Fix positive integers $t$ and $N$, and write $[N]=\{1,\ldots,N\}$.
Consider a subset $A\subseteq[N]$ with the following simultaneous
avoidance property: for every pair of integers $1\le m,n\le t$,
the equation
\[
                              mx+ny=(m+n)z
\]
has no nontrivial solution with $x,y,z\in A$. Here a solution is
trivial exactly when $x=y=z$. Since $m,n$ are positive, any solution
with two equal variables is already trivial; a prohibited solution
therefore uses three distinct elements. The coefficients are ordinary
integers, and equality is in $\mathbb Z$, without reduction modulo
any number.

Define
\[
 F_t(N)=\max\{|A|:A\subseteq[N]\text{ satisfies all the above
                                          avoidance conditions}\}.
\]
The problem is to determine how large $F_t(N)$ can be, particularly
its asymptotic growth as $N\to\infty$ with $t$ fixed and the dependence
of useful estimates on $t$.
The same set $A$ must avoid all $t^2$ equations at once; choosing
a different set for each coefficient pair is not sufficient.
For $m=n=1$, the excluded configuration is an ordinary three-term
arithmetic progression, while other coefficient pairs forbid further
rational weighted averages between two selected integers. Pairs of
proportional coefficients can specify the same equation, but retaining
them does not affect the definition. Both the ambient interval and
the cardinality objective are part of the extremal question.
\subsection{Short English statement}
How large can an interval subset be if it simultaneously avoids every nontrivial three-variable weighted-average equation with coefficients at most t?
\subsection{Sources}
\begin{itemize}
\item[S1] ULAM.AI and the UnsolvedMath Contributors, supplied problems.json, record 20001112 (AIM-COMBINATORICS-0237), AIM Workshop Problem Lists. Catalogue identity and source transcription; CC BY 4.0.
\url{https://huggingface.co/datasets/ulamai/UnsolvedMath}.
\item[S2] American Institute of Mathematics, Recent trends in additive combinatorics, item 2.12. Original workshop question underlying AIM-COMBINATORICS-0237.
\url{https://aimath.org/WWN/additivecomb/additivecomb.pdf}.
\end{itemize}
\end{document}
